\section*{Розділ III. Повноваження ЦВКс}
\addcontentsline{toc}{section}{Розділ III. Повноваження ЦВКс}

\subsection*{3.1. Загальні повноваження}
\addcontentsline{toc}{subsection}{3.1. Загальні повноваження}
    До повноважень ЦВКс належить:
        \begin{enumerate}[label=\alph*)]
            \item Організація та проведення виборів до органів студентського самоврядування відповідно до Положення про ОСС та цього Положення;
            \item Забезпечення реалізації виборчих прав студентів Університету;
            \item Контроль за дотриманням виборчого законодавства всіма учасниками виборчого процесу;
            \item Розгляд скарг та заяв з питань, пов'язаних з виборчим процесом;
            \item Прийняття рішень з питань, віднесених до її компетенції цим Положенням.
        \end{enumerate}

\subsection*{3.2. Повноваження щодо організації виборів}
\addcontentsline{toc}{subsection}{3.2. Повноваження щодо організації виборів}
    3.2.1. ЦВКс організовує та проводить:

        \begin{enumerate}[label=\alph*)]
            \item Прямі вибори Голови Студентської ради Київського авіаційного інституту (СР КАІ);
            \item Прямі вибори Голови Студентської ради студмістечка (СР СМ) (у разі її створення);
            \item Прямі вибори Голів студентських рад факультетів/інститутів (СРФ/СРІ);
            \item Вибори делегатів на Конференцію студентів Університету (КСУ) від факультетів/інститутів та гуртожитків;
            \item Вибори представників студентів до Вченої ради Університету та вчених рад факультетів/інститутів.
        \end{enumerate}

    3.2.2. ЦВКс не організовує внутрішні вибори керівних органів студентського самоврядування (обрання заступників, секретарів, членів президій тощо), які проводяться самими органами відповідно до їх внутрішніх положень.

    3.2.3. ЦВКс може надавати організаційно-технічну підтримку виборів членів СУД за дорученням КСУ відповідно до процедур, визначених Регламентом КСУ, не втручаючись у процедури внутрішнього обрання СУД.

    3.2.4. ЦВКс має право створювати дільничні виборчі комісії за потребою для забезпечення організації голосування, зокрема:

        \begin{enumerate}[label=\alph*)]
            \item При кількості виборців понад 1000 осіб;
            \item При проведенні виборів у віддалених підрозділах Університету;
            \item При технічній необхідності розосередження виборчих дільниць.
        \end{enumerate}

    3.2.5. Дільничні виборчі комісії створюються у складі 3-5 осіб і діють під керівництвом та контролем ЦВКс з дотриманням принципів неупередженості та професійності.

\subsection*{3.3. Нормотворчі повноваження}
\addcontentsline{toc}{subsection}{3.3. Нормотворчі повноваження}
    3.3.1. ЦВКс має право:

        \begin{enumerate}[label=\alph*)]
            \item Розробляти та подавати на затвердження КСУ загальні принципи та процедури проведення виборів до органів студентського самоврядування;
            \item Приймати технічні регламенти проведення виборів, що деталізують процедури, затверджені КСУ;
            \item Розробляти та подавати на затвердження КСУ технічні регламенти електронного голосування, включаючи визначення платформ, процедур безпеки, алгоритмів ідентифікації та забезпечення таємності волевиявлення;
            \item Затверджувати форми виборчих документів (бюлетенів, протоколів, заяв тощо);
            \item Визначати технічні аспекти виборчого процесу (способи голосування, порядок підрахунку голосів, вимоги до виборчих дільниць тощо).
        \end{enumerate}

    3.3.2. Технічні регламенти та інші нормативні акти ЦВКс не можуть суперечити Положенню про ОСС, цьому Положенню та рішенням КСУ.

    3.3.3. Технічні регламенти електронного голосування, затверджені КСУ, застосовуються при проведенні всіх виборів, включаючи вибори членів ЦВКс.

    3.3.4. Проєкти нормативних актів ЦВКс, що стосуються загальних принципів виборчого процесу, підлягають попередньому обговоренню з представниками студентської спільноти.

\subsection*{3.4. Повноваження щодо реєстрації кандидатів}
\addcontentsline{toc}{subsection}{3.4. Повноваження щодо реєстрації кандидатів}
    3.4.1. ЦВКс:

        \begin{enumerate}[label=\alph*)]
            \item Реєструє кандидатів на виборні посади в органах студентського самоврядування;
            \item Перевіряє відповідність кандидатів встановленим вимогам;
            \item Визначає конкретний порядок подання документів кандидатами та форми підтвердження їх статусу студента;
            \item Приймає рішення про реєстрацію або відмову в реєстрації кандидатів;
            \item Скасовує реєстрацію кандидатів у випадках, передбачених виборчим законодавством.
        \end{enumerate}

    3.4.2. Рішення ЦВКс про відмову в реєстрації або скасування реєстрації кандидата приймається більшістю голосів від загального складу ЦВКс з обов'язковим обґрунтуванням.

\subsection*{3.5. Контрольні та санкційні повноваження}
\addcontentsline{toc}{subsection}{3.5. Контрольні та санкційні повноваження}
    3.5.1. За порушення виборчого законодавства ЦВКс має право застосовувати такі санкції:

        \begin{enumerate}[label=\alph*)]
            \item Винесення попередження учасникам виборчого процесу;
            \item Вимога усунути виявлені порушення у встановлений строк;
            \item Скасування реєстрації кандидатів за систематичні або грубі порушення;
            \item Визнання результатів голосування на окремих виборчих дільницях недійсними;
            \item Прийняття рішення про повторне проведення виборів у разі масштабних порушень.
        \end{enumerate}

    3.5.2. У разі виявлення порушень, які можуть становити підставу для дисциплінарної відповідальності, ЦВКс має право звернутися до Студентської уповноваженої делегації (СУД) з відповідним поданням.

    3.5.3. Рішення ЦВКс про застосування санкцій можуть бути оскаржені до КСУ протягом 3 робочих днів з дня їх прийняття.

\subsection*{3.6. Інформаційні повноваження}
\addcontentsline{toc}{subsection}{3.6. Інформаційні повноваження}
    3.6.1. ЦВКс:

        \begin{enumerate}[label=\alph*)]
            \item Оприлюднює інформацію про призначення виборів, порядок та строки їх проведення;
            \item Забезпечує інформування студентів про кандидатів та їх програми;
            \item Оголошує результати виборів та оприлюднює підсумкові протоколи;
            \item Веде офіційну статистику проведених виборів.
        \end{enumerate}

    3.6.2. Інформація про виборчий процес оприлюднюється на офіційних інформаційних ресурсах органів студентського самоврядування та Університету.

\subsection*{3.7. Повноваження щодо встановлення результатів виборів}
\addcontentsline{toc}{subsection}{3.7. Повноваження щодо встановлення результатів виборів}
    3.7.1. ЦВКс:

        \begin{enumerate}[label=\alph*)]
            \item Здійснює підрахунок голосів та встановлює результати виборів;
            \item Приймає рішення про визнання виборів дійсними або недійсними;
            \item Складає підсумкові протоколи про результати виборів;
            \item Передає документи про результати виборів відповідним органам для затвердження (КСУ -- для результатів виборів Голів СР КАІ та СР СМ).
        \end{enumerate}

    3.7.2. Рішення ЦВКс про результати виборів приймається не пізніше ніж через 3 робочі дні після завершення голосування.